\documentclass[aspectratio=169, 10pt]{beamer}
\usepackage[fontset=fandol]{ctex}
\usepackage{amsmath, amssymb, mathtools}
\usepackage{booktabs}
\usepackage{graphicx}
\usepackage{xcolor}
\usepackage{tikz}
\usetikzlibrary{arrows.meta, positioning, shapes.geometric}

% -- Theme ---------------------------------------------------------------
\usetheme{Madrid}
\usecolortheme{dolphin}
\usefonttheme{professionalfonts}
\setbeamertemplate{navigation symbols}{}

% -- Compact spacing ------------------------------------------------------
\setbeamersize{text margin left=6mm, text margin right=6mm}
\renewcommand{\baselinestretch}{0.92}
\setlength{\parskip}{0pt}
\setbeamertemplate{itemize items}[circle]
\setbeamerfont{block title}{size=\small,series=\bfseries}
\setbeamerfont{block body}{size=\small}
\setbeamerfont{block title alerted}{size=\small,series=\bfseries}
\setbeamerfont{block body alerted}{size=\small}

% -- Colors --------------------------------------------------------------
\definecolor{myblue}{HTML}{2B5EA7}
\definecolor{accentgreen}{HTML}{2E8B57}
\definecolor{accentred}{HTML}{C0392B}
\definecolor{accentorange}{HTML}{D35400}
\definecolor{lightbg}{HTML}{F5F7FA}
\setbeamercolor{frametitle}{bg=myblue, fg=white}
\setbeamercolor{structure}{fg=myblue}
\setbeamercolor{block title}{bg=myblue!85, fg=white}
\setbeamercolor{block body}{bg=myblue!8}
\setbeamercolor{alerted text}{fg=accentred}
\setbeamercolor{block title alerted}{bg=accentred, fg=white}
\setbeamercolor{block body alerted}{bg=accentred!10}

% -- Auto section title page ---------------------------------------------
\AtBeginSection[]{
  \begin{frame}
    \vfill\centering
    \usebeamerfont{title}\textcolor{myblue}{\insertsectionhead}\par
    \vfill
  \end{frame}
}

% -- Footer with GitHub URL ----------------------------------------------
\setbeamertemplate{footline}{%
  \leavevmode%
  \hbox{%
  \begin{beamercolorbox}[wd=.55\paperwidth,ht=2.5ex,dp=1ex,leftskip=2ex]{author in head/foot}%
    \tiny\texttt{github.com/dechang64/organoid-fl}%
  \end{beamercolorbox}%
  \begin{beamercolorbox}[wd=.35\paperwidth,ht=2.5ex,dp=1ex,center]{author in head/foot}%
    \tiny\insertshorttitle%
  \end{beamercolorbox}%
  \begin{beamercolorbox}[wd=.10\paperwidth,ht=2.5ex,dp=1ex,center]{author in head/foot}%
    \tiny\insertframenumber{} / \inserttotalframenumber%
  \end{beamercolorbox}}%
  \vskip0pt%
}

% -- Metadata ------------------------------------------------------------
\title[Organoid-FL]{Organoid-FL: 类器官检测联邦学习平台}
\subtitle{项目进展汇报 v5}
\date{2026年7月}

\begin{document}

% ========================================================================
\begin{frame}[plain]
  \titlepage
\end{frame}

\begin{frame}{目录}
  \tableofcontents
\end{frame}

% ========================================================================
\section{问题与背景}

\begin{frame}{什么是类器官？}
  \begin{columns}[T]
    \column{0.55\textwidth}
    \begin{block}{类器官（Organoid）}
      \begin{itemize}
        \item 体外培养的微型器官（肺、肠、脑等）
        \item 用于药物筛选、个性化医疗、疾病建模
        \item 显微镜下成像后需\textbf{人工计数和形态分析}
        \item 通量低、主观性强、跨实验室一致性差
      \end{itemize}
    \end{block}
    \vspace{0em}
    \begin{alertblock}{核心痛点}
      类器官图像分析依赖人工，成为药物研发管线的瓶颈。
      自动化检测是解决方案，但面临多个技术挑战。
    \end{alertblock}
    \column{0.40\textwidth}
    \begin{block}{技术挑战}
      \begin{itemize}
        \item \alert{超大图}：6K$\times$5.7K 像素
        \item \alert{密集小目标}：单图数百个 organoid
        \item \alert{标注噪声}：多标注者不一致
        \item \alert{数据分散}：多中心数据无法共享
      \end{itemize}
    \end{block}
  \end{columns}
\end{frame}

\begin{frame}{为什么需要联邦学习？}
  \begin{block}{数据隐私与数据分散的矛盾}
    \begin{itemize}
      \item 多家医院/实验室各自有类器官图像数据
      \item 涉及患者来源 → \alert{无法集中共享}
      \item 单中心数据量不足 → 模型泛化性差
      \item 各中心培养条件/成像设备不同 → 数据分布异构（non-IID）
    \end{itemize}
  \end{block}
  \vspace{0em}
  \begin{block}{联邦学习（Federated Learning）}
    \begin{itemize}
      \item \textbf{数据不动，模型动}：各中心本地训练，只上传模型参数
      \item \textbf{数据不动，知识动}：聚合后的全局模型吸收所有中心的知识
      \item \textbf{数据不动，推理动}：全局模型下发给各中心本地推理
    \end{itemize}
  \end{block}
  \vspace{0em}
  \centering\alert{FL + Organoid 是未被探索的交叉方向}
\end{frame}

% ========================================================================
\section{踩坑总结}

\begin{frame}{踩坑总结：数据集层面}
  \begin{alertblock}{铁律：使用任何新数据集前，必须先查官方文档}
    \begin{itemize}
      \item MultiOrg 标注 JSON 使用 napari 的 \textbf{(row, col) = (y, x)} 坐标系，不是常规 (x, y)
      \item 论文用\textbf{单类} organoid（Normal/Macros 只是培养条件，不是检测类别）
      \item 论文用 \textbf{512 patch + 丢弃边界 bbox + 全图滑动窗口推理}，不是 640 patch 上直接评估
    \end{itemize}
  \end{alertblock}
  \vspace{0em}
  \begin{block}{三个根本性错误 → 100 epoch mAP=0}
    \begin{tabular}{p{3.5cm}p{4cm}p{4cm}}
      \toprule
      & \textbf{错误做法} & \textbf{正确做法（论文方法）} \\
      \midrule
      类别 & 2类（Normal/Macros） & \textbf{单类} organoid \\
      边界 bbox & 裁剪到 patch 边缘 & \textbf{丢弃}边界 bbox \\
      评估 & patch 上直接评估 & 全图滑动窗口推理 \\
      \bottomrule
    \end{tabular}
  \end{block}
  \vspace{0em}
  \begin{block}{坐标系错误 → 标注越界}
    napari [y,x] 格式被当 [x,y] 读 → 标注 y 最大 5845 $>$ 图片高 5720，交换后 x 最大 5845 $<$ 宽 6385 ✓
  \end{block}
\end{frame}

\begin{frame}{踩坑总结：代码层面}
  \begin{block}{2026-06-28 系统审计发现 7 个 bug}
    \begin{tabular}{p{3.2cm}p{2.5cm}p{4.5cm}p{1.2cm}}
      \toprule
      \textbf{文件} & \textbf{Bug} & \textbf{影响} & \textbf{严重性} \\
      \midrule
      multiorg\_sam2.py & mask\_iou\_cropped offset x/y 交换 & 所有 mask mAP 数值错误 & \alert{Critical} \\
      multiorg\_sam2.py & mask\_iou() 死代码 & 无 & Minor \\
      multiorg\_sam2.py & load\_sam2 临时文件不清理 & 磁盘泄漏 & Medium \\
      multiorg\_tiling\_v3.py & PIL I;16→numpy segfault & Macros 类崩溃 & \alert{Critical} \\
      prepare\_sam2\_pseudo.py & cache return 跳过所有后续图 & 数据缺失 & \alert{Critical} \\
      prepare\_sam2\_pseudo.py & fallback mask 全图分配 & OOM 风险 & Medium \\
      prepare\_sam2\_pseudo.py & instance 编号错位 & 训练读错 mask & \alert{Critical} \\
      \bottomrule
    \end{tabular}
  \end{block}
  \vspace{0em}
  \begin{alertblock}{教训}
    \begin{itemize}
      \item 写完脚本必须端到端验证，不能只过 py\_compile
      \item offset/坐标格式必须统一标注 (x,y) 还是 (y,x)
      \item docstring 里的路径不能编造，必须从训练日志确认
    \end{itemize}
  \end{alertblock}
\end{frame}

\begin{frame}{踩坑总结：实验设计层面}
  \begin{block}{val 集 drop\_boundary 导致 mAP 暴跌 15pp}
    \begin{itemize}
      \item train 用 drop\_boundary=True（干净信号）✓
      \item val 也用 drop\_boundary=True → \alert{把模型正确检测变 false positive}
      \item 正确做法：val 必须 clip 到 patch 边缘（保留所有 GT）
    \end{itemize}
  \end{block}
  \vspace{0em}
  \begin{block}{GT 标注质量决定微调策略}
    \begin{tabular}{p{2.5cm}p{4cm}p{4cm}}
      \toprule
      & \textbf{鼠肝} & \textbf{MultiOrg} \\
      \midrule
      标注类型 & 红色折线（精确轮廓） & 4点多边形（位置标注） \\
      SAM2 微调 & \textbf{有效}（F1=93.9\%） & \textbf{有害}（mask mAP -4pp） \\
      正确做法 & mask supervision & 位置 prompt only \\
      \bottomrule
    \end{tabular}
  \end{block}
  \vspace{0em}
  \begin{alertblock}{核心教训}
    用粗糙 GT 微调精确 SAM2 = 负优化。GT 质量决定 supervision 方式。
  \end{alertblock}
\end{frame}

% ========================================================================
\section{MultiOrg 实验结果}

\begin{frame}{标注者选择对评估的影响}
  \begin{block}{同一模型，不同标注者评估}
    \begin{tabular}{lcccc}
      \toprule
      标注者 & Instances & mAP@0.5 & mAP@0.5:0.95 & Recall \\
      \midrule
      t0（初始） & 6465 & 66.9\% & 32.6\% & 60.2\% \\
      t1\_A（二次） & 3950 & 70.5\% & 38.7\% & 64.8\% \\
      \textbf{t1\_B（最精准）} & 4112 & \textbf{78.9\%} & \textbf{53.9\%} & \textbf{72.3\%} \\
      \bottomrule
    \end{tabular}
  \end{block}
  \vspace{0em}
  \begin{alertblock}{关键发现}
    \begin{itemize}
      \item 标注者选择可造成 \textbf{12pp mAP@0.5 差异}（66.9\% vs 78.9\%）
      \item t0 噪声最大（6465 instances，含误标）；t1\_B 最精简（4112 instances）
      \item \textbf{评估方法论和模型本身一样重要}——噪声标注会严重低估模型能力
    \end{itemize}
  \end{alertblock}
\end{frame}

\begin{frame}{MultiOrg 检测精度 vs SOTA}
  \begin{block}{mAP@0.5 对比（t1\_b 标注）}
    \begin{tabular}{lcc}
      \toprule
      配置 & mAP@0.5 & vs SOTA \\
      \midrule
      Faster R-CNN（论文） & 68.4\% & $-$5.5pp \\
      YOLOv3（论文） & 70.3\% & $-$3.6pp \\
      RTMDet（论文） & 63.2\% & $-$10.7pp \\
      \textbf{SSD（论文 SOTA）} & \textbf{73.9\%} & — \\
      \midrule
      YOLOv12s patch 级 & 78.1\% & +4.2pp \\
      RF-DETR SAHI+Soft-NMS & 77.15\% & +3.3pp \\
      \textbf{RF-DETR ensemble baseline} & \textbf{77.8\%} & \textbf{+3.9pp} \\
      \bottomrule
    \end{tabular}
  \end{block}
  \vspace{0em}
  \begin{alertblock}{结论}
    \begin{itemize}
      \item \textbf{我们已是 MultiOrg SOTA}，超论文最优 SSD +3.9pp
      \item 论文 2024.10 发表，无后续工作超过我们
      \item 距 80\% 门槛差 2.2pp，瓶颈是数据量（356张 vs Orga-Dete 4008张）
    \end{itemize}
  \end{alertblock}
\end{frame}

\begin{frame}{SAM2 分割实验：自蒸馏方案}
  \begin{block}{问题：4点多边形 GT 太粗糙}
    MultiOrg 标注只标 4 个点（位置），不是精确轮廓。直接用 4点多边形微调 SAM2 = 负优化（mask mAP -4pp）。
  \end{block}
  \vspace{0em}
  \begin{block}{方案 A：自蒸馏（Self-Distillation）}
    \begin{enumerate}
      \item GT 4点多边形 → bbox（位置 prompt）
      \item SAM2 zero-shot 用 bbox 做 prompt → 伪 mask（精确轮廓）
      \item 伪 mask 作为微调 supervision（替代粗糙的 4点多边形）
    \end{enumerate}
  \end{block}
  \vspace{0em}
  \begin{block}{训练结果（5 epoch, 3060 12GB）}
    \begin{tabular}{ccccc}
      \toprule
      Epoch & Train IoU & Val IoU & Val Dice & 耗时 \\
      \midrule
      1 & 0.8437 & 0.8419 & 0.9074 & 70min \\
      3 & 0.8646 & 0.8594 & 0.9181 & 70min \\
      \textbf{5} & \textbf{0.8814} & \textbf{0.8720} & \textbf{0.9255} & 70min \\
      \bottomrule
    \end{tabular}
  \end{block}
\end{frame}

\begin{frame}{SAM2 对照实验结果}
  \begin{block}{全量 55 张测试，mask IoU bug 修复后}
    \begin{tabular}{lcccc}
      \toprule
      \textbf{配置} & \textbf{bbox mAP50} & \textbf{mask mAP50} & \textbf{F1} & \textbf{FP} \\
      \midrule
      Zero-shot SAM2 & 77.15\% & \textbf{76.39\%} & 43.9\% & 11553 \\
      Finetuned v2 (4点GT) & 77.15\% & 75.98\% & 43.9\% & 11553 \\
      Finetuned pseudo & 77.13\% & 75.89\% & 43.9\% & 11553 \\
      \bottomrule
    \end{tabular}
  \end{block}
  \vspace{0em}
  \begin{alertblock}{结论：微调 = 负优化}
    \begin{itemize}
      \item bbox mAP50 三者一致（SAM2 不改检测框）
      \item mask mAP50：zero-shot 76.39\% $>$ finetuned 75.98\% $>$ pseudo 75.89\%
      \item 根因：4 点粗糙 GT 微调精确 SAM2 = 负优化
      \item 自蒸馏天花板 = zero-shot 本身质量
    \end{itemize}
  \end{alertblock}
  \vspace{0em}
  \begin{block}{FP 抑制全面失败}
    \begin{tabular}{lc}
      \toprule
      方案 & 结论 \\
      \midrule
      形态学过滤 (circ/solid/ar) & 砍 FP $<$ 2\%，无效 \\
      DINOv2+DPMM 二级验证 & PR-AUC=0.29 $<$ 随机 0.5，完全不可分 \\
      HNM (Hard Negative Mining) & 灾难性遗忘，mAP 暴跌 40pp \\
      \bottomrule
    \end{tabular}
  \end{block}
\end{frame}

% ========================================================================
\section{鼠肝实验规划}

\begin{frame}{鼠肝实验：已有基础}
  \begin{columns}[T]
    \column{0.48\textwidth}
    \begin{block}{数据}
      \begin{itemize}
        \item Batch 1: 10对图（2592$\times$1944, 4X），23个类器官
        \item Batch 2: 10对图（4000$\times$3000, 不同相机），41个类器官
        \item \textbf{红色折线轮廓标注}（精确轮廓，非4点）
        \item OpenCV 提取轮廓 → YOLO 格式 bbox
      \end{itemize}
    \end{block}
    \column{0.48\textwidth}
    \begin{block}{Batch 1 结果}
      \begin{itemize}
        \item RF-DETR few-shot（8张, 1ep）
        \item + zero-shot SAM2 分割
        \item \textbf{F1 = 93.9\%}
        \item SAM2 mask 和人工红线格式一致
      \end{itemize}
    \end{block}
  \end{columns}
  \vspace{0em}
  \begin{alertblock}{Batch 2 域偏移测试}
    Batch 1 模型 → Batch 2 推理 → F1=0\%（完全失败）。根因：不同相机/光照/对比度，不是分辨率问题。
  \end{alertblock}
\end{frame}

\begin{frame}{鼠肝实验：为什么值得做？}
  \begin{block}{鼠肝 vs MultiOrg 对比}
    \begin{tabular}{p{3.5cm}p{4.5cm}p{4cm}}
      \toprule
      & \textbf{鼠肝} & \textbf{MultiOrg} \\
      \midrule
      标注质量 & 红色折线（精确轮廓） & 4点多边形（位置） \\
      SAM2 微调 & \textbf{有效}（F1=93.9\%） & 有害（-4pp） \\
      形态学过滤 & 有效（conf+circ） & 无效 \\
      数据量 & 20张（few-shot） & 411张 \\
      域偏移 & \alert{有}（batch 1→2 F1=0） & 无 \\
      \bottomrule
    \end{tabular}
  \end{block}
  \vspace{0em}
  \begin{block}{鼠肝实验的独特价值}
    \begin{itemize}
      \item \textbf{精确 GT}：折线轮廓 → SAM2 微调有效 → 可以做 mask 级评估
      \item \textbf{域偏移}：不同实验室/相机/光照 → FL 的天然场景
      \item \textbf{Few-shot}：数据少 → FL 价值更大（单中心不够训）
      \item \textbf{已有 pipeline}：RF-DETR + SAM2 + 形态学过滤已验证
    \end{itemize}
  \end{block}
\end{frame}

\begin{frame}{鼠肝实验：三阶段计划}
  \begin{block}{Phase 1：单中心精度验证（1-2周）}
    \begin{enumerate}
      \item Batch 1（10张）训练 RF-DETR few-shot + SAM2 微调（已有 F1=93.9\%）
      \item Batch 2（10张）作为独立 test set → 模拟跨域泛化
      \item 对比 zero-shot SAM2 vs finetuned SAM2 的 mask 质量
      \item \textbf{关键验证}：精确 GT 微调 SAM2 是否真有增益（vs MultiOrg 的负优化）
    \end{enumerate}
  \end{block}
  \vspace{0em}
  \begin{block}{Phase 2：跨域泛化 + 图像预处理标准化（2-3周）}
    \begin{enumerate}
      \item Batch 1 → Batch 2 域偏移量化（当前 F1=0\%）
      \item 图像预处理标准化模块：
        \begin{itemize}
          \item 分辨率归一化（到固定 μm/px）
          \item 对比度归一化（CLAHE）
          \item 背景扣除
          \item 位深统一（16bit→8bit）
        \end{itemize}
      \item 验证预处理后跨域 F1 是否提升
    \end{enumerate}
  \end{block}
\end{frame}

\begin{frame}{鼠肝实验：FL 场景设计}
  \begin{block}{Phase 3：联邦学习实验（3-4周）}
    \begin{enumerate}
      \item \textbf{Client 划分}：
        \begin{itemize}
          \item Client A: Batch 1（2592$\times$1944, 4X 相机）
          \item Client B: Batch 2（4000$\times$3000, 不同相机）
          \item Client C: 数据增强生成（模拟第三中心）
        \end{itemize}
      \item \textbf{Non-IID 设计}：
        \begin{itemize}
          \item 分辨率异质（不同 μm/px）
          \item 光照异质（不同对比度/亮度）
          \item 标注密度异质（不同 organoid 数量/图）
        \end{itemize}
      \item \textbf{聚合策略对比}：FedAvg vs EWA vs FedProx
      \item \textbf{评估}：mask mAP50 + bbox mAP50 + 跨域 F1
    \end{enumerate}
  \end{block}
  \vspace{0em}
  \begin{alertblock}{预期贡献}
    \begin{itemize}
      \item 首个 FL + 类器官检测的跨域泛化研究
      \item 图像预处理标准化模块（多实验室平台的刚需）
      \item EWA Effective Interval 在跨模态场景的验证
    \end{itemize}
  \end{alertblock}
\end{frame}

\begin{frame}{鼠肝实验：技术路线图}
  \centering
  \begin{tikzpicture}[
    node distance=0.6cm and 1.0cm,
    box/.style={draw, rounded corners, minimum width=3.2cm, minimum height=0.8cm,
                font=\small, align=center, fill=myblue!10},
    phase/.style={draw, rounded corners, minimum width=3.5cm, minimum height=0.7cm,
                  font=\small\bfseries, align=center, fill=myblue!85, text=white},
    arrow/.style={-Stealth, thick, myblue}
  ]
    \node[phase] (p1) {Phase 1: 单中心精度};
    \node[phase, right=of p1] (p2) {Phase 2: 跨域泛化};
    \node[phase, right=of p2] (p3) {Phase 3: 联邦学习};

    \node[box, below=of p1] (b1) {RF-DETR few-shot\\+ SAM2 微调};
    \node[box, below=of p2] (b2) {预处理标准化\\CLAHE + 分辨率归一};
    \node[box, below=of p3] (b3) {3-client FL\\FedAvg vs EWA};

    \node[box, below=of b1, fill=accentgreen!15] (r1) {F1=93.9\% ✓};
    \node[box, below=of b2, fill=accentorange!15] (r2) {目标: F1>80\%};
    \node[box, below=of b3, fill=accentorange!15] (r3) {目标: 超单中心};

    \draw[arrow] (p1) -- (p2);
    \draw[arrow] (p2) -- (p3);
    \draw[arrow] (b1) -- (r1);
    \draw[arrow] (b2) -- (r2);
    \draw[arrow] (b3) -- (r3);
  \end{tikzpicture}
  \vspace{0.5em}
  \begin{itemize}
    \item Phase 1 已有基础（F1=93.9\%），重点验证 SAM2 微调增益
    \item Phase 2 解决域偏移（batch 1→2 F1=0\%），预处理标准化是关键
    \item Phase 3 FL 实验，鼠肝的精确 GT 让 mask 级评估有意义
  \end{itemize}
\end{frame}

% ========================================================================
\section{联邦学习}

\begin{frame}{EWA 熵加权聚合策略}
  \begin{block}{核心思想}
    传统 FedAvg 按样本量均匀加权，忽略客户端模型质量差异。\\
    EWA 用验证集熵作为信号，\textbf{给更好的客户端更高权重}：
    \[
      w_c = \frac{\exp(-H_c / \alpha)}{\sum_{c'} \exp(-H_{c'} / \alpha)}
    \]
    其中 $H_c$ 是客户端 $c$ 在验证集上的预测熵，$\alpha$ 是温度参数。
  \end{block}
  \vspace{0em}
  \begin{block}{EWA Effective Interval 理论}
    \begin{itemize}
      \item \textbf{IID}：质量差异小 → EWA 退化 → \alert{无优势}
      \item \textbf{适度异质}：差异可区分 → EWA 最优 → \alert{+1.4\textasciitilde 1.5pp}
      \item \textbf{极端异质}：信号噪声大 → \alert{+1.1pp}
    \end{itemize}
  \end{block}
\end{frame}

% ========================================================================
\section{总结}

\begin{frame}{Ensemble 实验：RF-DETR + YOLOv12}
  \begin{block}{三种融合策略结果}
    \begin{tabular}{lccc}
      \toprule
      策略 & mAP50 & mAP50-95 & vs RF-DETR \\
      \midrule
      RF-DETR 单模型 & 77.8\% & 48.8\% & — \\
      YOLOv12s 单模型 & 62.3\% & 37.2\% & $-$15.5pp \\
      \midrule
      WBF（SOTA 融合） & 63.2\% & 38.3\% & $-$14.5pp \\
      Intersection & 63.1\% & 39.7\% & $-$14.7pp \\
      Union & 70.0\% & 42.8\% & $-$7.8pp \\
      \bottomrule
    \end{tabular}
  \end{block}
  \begin{alertblock}{失败原因：YOLOv12s 是严格劣势模型}
    \begin{itemize}
      \item YOLO recall 80.4\% $\ll$ RF-DETR 96.2\%，无独有 TP
      \item Ensemble TP=4745 $<$ RF-DETR TP=4767 → YOLO 未补充任何 recall
      \item WBF 的 T/N 机制反而拉低 RF-DETR 高分 TP 的 score
      \item \textbf{结论：ensemble 需要两个精度接近的模型才能受益}
    \end{itemize}
  \end{alertblock}
\end{frame}

\begin{frame}{全面文献调研：MultiOrg SOTA}
  \begin{block}{MultiOrg 数据集 SOTA（mAP@0.5, t1\_b）}
    \begin{tabular}{llc}
      \toprule
      方法 & 来源 & mAP50 \\
      \midrule
      SSD & MultiOrg 论文 SOTA & 73.9\% \\
      YOLOv3 & MultiOrg 论文 & 70.3\% \\
      Faster R-CNN & MultiOrg 论文 & 68.4\% \\
      \textbf{RF-DETR（我们）} & \textbf{本研究} & \textbf{77.8\%} \\
      \bottomrule
    \end{tabular}
  \end{block}
  \begin{block}{其他 organoid 数据集 SOTA}
    \begin{tabular}{llcl}
      \toprule
      方法 & mAP50 & 数据集 & 数据量 \\
      \midrule
      Deliod & 87.5\% & Intestinal & — \\
      Orga-Dete & 81.4\% & Lung & 4008 张 \\
      Orga-Dete & 92.5\% & Duodenal & — \\
      Orga-Dete & 84.4\% & Murine intestinal & — \\
      OrgLine & N/A (PQ) & Multi-source & — \\
      \bottomrule
    \end{tabular}
  \end{block}
\end{frame}

\begin{frame}{Orga-Dete 三模块迁移实验：结果}
  \begin{alertblock}{结论：三模块全面失败，停止该方向}
    \begin{tabular}{lccc}
      \toprule
      \textbf{配置} & \textbf{mAP50} & \textbf{vs baseline} & \textbf{结论} \\
      \midrule
      yolo12s baseline & 62.3\% & — & 参考 \\
      yolo11n + MPCA & 44.2\% & — & 容量瓶颈 \\
      yolo11n + MPCA + EMASlideLoss & 44.0\% & $-$0.2pp & 噪声级 \\
      yolo12s + MPCA (random init) & 43.5\% & $-$18.8pp & 破坏预训练 \\
      yolo12s + MPCA (identity init) & 43.9\% & $-$18.4pp & init 不是根因 \\
      yolo12s + MPCA (identity) + EMASlideLoss & 44.5\% & $-$17.8pp & 噪声级 \\
      \bottomrule
    \end{tabular}
  \end{alertblock}
  \vspace{0em}
  \begin{block}{MPCA Identity Init 修复（commit 0cd3169）}
    \begin{itemize}
      \item \textbf{问题}：Sigmoid 随机初始化 $\approx$ 0.5 $\rightarrow$ 特征方差减半 $\rightarrow$ 破坏预训练
      \item \textbf{修复}：weight=0, bias=5 $\rightarrow$ sigmoid(5)$\approx$1.0 $\rightarrow$ output$\approx$x（identity）
      \item \textbf{验证}：std ratio 0.50$\rightarrow$0.99，预训练匹配 48.6\%$\rightarrow$99.2\%
      \item \textbf{但 mAP 几乎不变}（+0.38pp）——init 不是根因
    \end{itemize}
  \end{block}
\end{frame}

\begin{frame}{Orga-Dete 失败根因分析}
  \begin{columns}[T]
    \column{0.48\textwidth}
    \begin{block}{模型容量 vs mAP 强相关}
      \begin{tabular}{lc}
        \toprule
        模型 & mAP50 \\
        \midrule
        RF-DETR small (32M) & 77.8\% \\
        YOLOv12s (9.4M) & 62.3\% \\
        YOLOv11n+MPCA (2.6M) & 44.2\% \\
        YOLOv12s+MPCA (9.6M) & 43.9\% \\
        \bottomrule
      \end{tabular}
      \vspace{0.3em}
      
      每 3.5x 参数 $\approx$ +18pp
    \end{block}
    \column{0.48\textwidth}
    \begin{block}{五个根因}
      \begin{enumerate}
        \item \textbf{模型容量瓶颈}：2.6M vs 32M
        \item \textbf{MPCA 架构不兼容}：C3k2 vs A2C2f
        \item \textbf{单类检测无增益}：坐标注意力对多类有用
        \item \textbf{数据集差异}：4008张 vs 411张
        \item \textbf{BiFPN 未验证}：YAML 不兼容
      \end{enumerate}
    \end{block}
  \end{columns}
  \vspace{0em}
  \begin{alertblock}{铁律：乘性 Attention 必须 Identity Init}
    任何加在预训练模型上的乘性 attention 模块（SE、CBAM、CA、MPCA），初始状态必须是 identity，否则 sigmoid$\approx$0.5 会破坏预训练特征。验证方法：output/input std ratio $\approx$ 1.0。
  \end{alertblock}
\end{frame}

\begin{frame}{总结}
  \begin{columns}[T]
    \column{0.48\textwidth}
    \begin{block}{已完成}
      \begin{itemize}
        \item[\checkmark] SOTA 全面调研（MultiOrg + Orga-Dete + OrgLine）
        \item[\checkmark] 数据预处理管线（tiling + 多标注者）
        \item[\checkmark] 标注 bug 修复（mAP +26pp）
        \item[\checkmark] \textbf{MultiOrg SOTA: 77.8\%（超论文 +3.9pp）}
        \item[\checkmark] SAM2 三轮实验：微调=负优化确认
        \item[\checkmark] FP 抑制全面调研：DINOv2/形态学/HNM 均无效
        \item[\checkmark] Ensemble 实验：WBF/Union/Intersection 全面失败
        \item[\checkmark] \textbf{Orga-Dete 三模块迁移：全面失败}
        \item[\checkmark] MPCA Identity Init 铁律发现
        \item[\checkmark] 鼠肝 batch 1 验证（F1=93.9\%）
      \end{itemize}
    \end{block}
    \column{0.48\textwidth}
    \begin{block}{下一步}
      \begin{itemize}
        \item[$\rightarrow$] \textbf{RF-DETR + SAHI 优化冲 80\%}
        \item[$\rightarrow$] yolo12m (20M) 容量升级评估
        \item[$\circ$] 鼠肝 Phase 1-3 实验
        \item[$\circ$] 图像预处理标准化模块
        \item[$\circ$] FL + Organoid 系统实验
        \item[$\circ$] 论文撰写 + 专利申请
      \end{itemize}
    \end{block}
  \end{columns}
\end{frame}

\begin{frame}[allowframebreaks]{参考文献}
  \footnotesize
  \begin{thebibliography}{99}
    \bibitem{bukas2024} Bukas et al. \textit{MultiOrg: A Multi-rater Organoid-detection Dataset}. NeurIPS, 2024.
    \bibitem{tian2025} Tian et al. \textit{YOLOv12: Attention-Centric Real-Time Object Detectors}. arXiv:2502.12524, 2025.
    \bibitem{solawetz2025} Solawetz et al. \textit{RF-DETR: Real-Time DEtection TRansformer}. ICLR, 2026.
    \bibitem{bodla2017} Bodla et al. \textit{Soft-NMS -- Improving Object Detection with One Line of Code}. ICCV, 2017.
    \bibitem{akyon2022} Akyon et al. \textit{Slicing Aided Hyper Inference (SAHI)}. CVPRW, 2022.
    \bibitem{bulten2021} Bulten et al. \textit{Overcoming the Limitations of Patch-based Learning to Detect Cancer in Whole Slide Images}. PMC, 2021.
    \bibitem{lu2025} Lu et al. \textit{Thinking with Visual Primitives}. DeepSeek-AI, 2025.
    \bibitem{kirillov2023} Kirillov et al. \textit{Segment Anything (SAM)}. ICCV, 2023.
    \bibitem{solovyev2021} Solovyev et al. \textit{Weighted Boxes Fusion: Ensembling boxes from different object detection models}. IVC, 2021.
    \bibitem{huang2025} Huang et al. \textit{Orga-Dete: An Improved Lightweight Deep Learning Model for Lung Organoid Detection}. Applied Sciences, 2025.
    \bibitem{deng2026} Deng et al. \textit{OrgLine: A versatile pipeline for organoid morphometry using detector-guided prompts}. Cell Reports Methods, 2026.
  \end{thebibliography}
\end{frame}

\begin{frame}[plain]
  \centering
  \vfill
  \Large\textcolor{myblue}{Thank You!}\\[0.8em]
  \normalsize Questions \& Discussion\\[1.5em]
  \footnotesize\texttt{github.com/dechang64/organoid-fl}
  \vfill
\end{frame}
\end{document}
